\subsection{Quy tắc chuẩn hóa và tham số thực thi}

Đầu vào cụ thể là \texttt{FleetPy\_Manhattan.zip} của bộ dữ liệu Manhattan,
phiên bản 1.0 trên Zenodo~\cite{fleetpydata2025}. Mỗi mẫu nguồn có một CSV
nhu cầu riêng; 108 yêu cầu là tập dẫn xuất tiếp theo, không phải toàn bộ
nhu cầu của khung giờ. Công cụ chuẩn hóa loại hàng sai định dạng, trùng
định danh, ngoài cửa sổ thời gian hoặc có hai đầu không thuộc cùng một
thành phần liên thông mạnh của mạng có hướng.

Việc chọn nút ưu tiên giữ được nhiều hàng nhu cầu. Công cụ bắt đầu từ cặp
điểm có nhiều hàng hợp lệ nhất, rồi thêm nút làm tăng số hàng có cả hai
đầu trong tập đã chọn. Khi hòa, nó dùng tần suất điểm, hạng HMAC và định
danh nút theo quy tắc cố định. Sau khi khóa tập 96 nút, các hàng được xếp
bằng HMAC-SHA256 từ băm kho nguồn, tên CSV và số thứ tự hàng; giữ 108 hàng
đầu tiên có cả hai đầu trong tập nút. Đây là lựa chọn tất định theo độ phủ
mạng và hạng băm, không phải mẫu ngẫu nhiên đơn của mọi chuyến taxi.
Các nhãn khóa và cửa sổ nằm trong \path{normalizer-configuration.json};
\path{selection-frame.json} giữ liên kết tới hàng nguồn.

Thời gian đến giữ khoảng cách trong CSV và được dịch về đầu ô; giây được
đổi sang mili giây bằng quy tắc làm tròn về số chẵn khi hòa. Mỗi cung trong
bao đóng dùng đường đi ngắn nhất trên mạng nguồn có hướng; đường đó có thể
qua nút ngoài tập 96 điểm. Không dùng đường ngược hoặc khoảng cách Euclid
để điền thời gian thiếu. Hệ số ban đầu lấy tại mốc 0 của tệp thời gian nguồn.

Đội xe ban đầu rỗng, không có khách đã nhận hoặc đang trên xe. Vị trí tám
xe là tám nút khác nhau đứng đầu hạng HMAC theo khóa chọn, băm kho nguồn,
định danh kịch bản và nút. Các nhánh của cùng ô dùng cùng đội xe; đổi kịch
bản hoặc cửa sổ có thể đổi vị trí. Bảng~\ref{bang:thamso} nêu các tham số
còn lại đã kiểm trực tiếp trong đầu vào và cấu hình bộ chạy.

\begin{table}[htbp]
\centering
\caption{Tham số chung của đầu vào và bộ chạy. Hạt giống và các cấu hình
trên cùng ô không được tính là mẫu độc lập.}
\label{bang:thamso}
\small
\begin{tabularx}{\linewidth}{@{}>{\raggedright\arraybackslash}p{4.1cm}>{\raggedright\arraybackslash}X@{}}
\toprule
\textbf{Tham số} & \textbf{Giá trị và quy tắc}\\
\midrule
Nhu cầu và xe & Mỗi yêu cầu một người; 108 yêu cầu, 8 xe bốn chỗ, xe ban đầu rỗng.\\
Cửa sổ đón & Từ lúc yêu cầu đến đến 600 giây sau đó; nhánh độ nhạy dùng 900 giây.\\
Thời gian trên xe tối đa & $1{,}5$ lần thời gian đi trực tiếp trong ảnh chụp ban đầu,
làm tròn về mili giây; không đặt lại theo từng cập nhật giao thông.\\
Phục vụ điểm dừng & Thời gian cố định lên và xuống xe bằng 0; vẫn có thể chờ đến giờ đón sớm nhất.\\
Cửa sổ và phần chạy nốt & Nhu cầu trong 7.200 giây, cho phép chạy nốt thêm 7.200 giây.
Bộ điều khiển có tối đa 24 mốc kiểm phần chạy nốt, cách nhau 600 giây.\\
Danh mục giữ lại & Tối đa 100 ứng viên mỗi xe, chèn tối đa hai yêu cầu mới;
giữ danh mục theo hồ sơ \texttt{service-set-stability-portfolio-v1}.\\
Ngân sách công việc & Sinh và kiểm: 10.000 đơn vị mỗi phần; bộ giải: 100.000 đơn vị,
thời gian tất định 1.000.000 micro giây. Đây không phải thời gian đáp ứng đã đo.\\
Bộ giải và hạt giống & OR-Tools 9.15.6755; hạt giống bộ chạy và bộ giải bằng 7;
chế độ liệt kê chính xác cho bài nhỏ tắt.\\
\bottomrule
\end{tabularx}
\end{table}

\subsection{Phiên bản và cấu hình theo lớp bằng chứng}

Bảng~\ref{bang:phienban} tách các quy tắc thực thi có thể làm đổi kết quả.
Các tổng H6 ở đây được đọc từ cây chạy lại \texttt{confirmatory-h6-panela-v2}.
Thiết kế và tiêu chí H6 gắn với biên nhận v5; v6 sửa ánh xạ đoạn đường xe
đang đi và cách băm nguồn bộ chuyển đổi rồi chạy lại thiết kế. Hai đợt dùng
cùng DLL bộ thực thi, nhưng bộ chuyển đổi và mốc mã ghi trong biên bản khác
nhau. Chạy lại không tạo một mẫu xác nhận độc lập hoặc thay tiêu chí H6.

\begin{table}[htbp]
\centering
\caption{Khác biệt giữa các lớp bằng chứng. Mã băm rút gọn giúp tra cứu;
định danh đầy đủ nằm trong biên nhận và bản kê của từng lần chạy.}
\label{bang:phienban}
\small
\begin{tabularx}{\linewidth}{@{}>{\raggedright\arraybackslash}p{2.1cm}>{\raggedright\arraybackslash}X@{}}
\toprule
\textbf{Lớp} & \textbf{Bộ thực thi, thế giới và quy tắc phục hồi}\\
\midrule
H6 & DLL \texttt{8445f053}; biên nhận v5 \texttt{84f6eff3}, v6 \texttt{e7bc41be}.
Ảnh chụp thời gian đi tĩnh; B1 dùng hồ sơ chi phí, C1 dùng vector cứng.
Chưa có phục hồi không tệ hơn tham chiếu.\\[4pt]
Tầng 2 & DLL \texttt{8b00dfbf}; bộ chạy hạn chót riêng.
Hệ số theo ngày được áp tại kích hoạt quyết định đầu tiên ở hoặc sau mỗi mốc 900 giây.
Không bật phục hồi không tệ hơn hoặc chế độ ghi nhận đón muộn.\\[4pt]
T38 & DLL \texttt{a631038b}; bộ chạy \texttt{runner-tier3-v2}, cấu hình WP4
\texttt{66e9b592}. Thế giới dừng đúng mỗi mốc 900 giây; số mũ hệ số bằng 1,
không thêm nhiễu hay sốc. Cả năm nhánh dùng phục hồi không tệ hơn tham chiếu
và chế độ ghi nhận đón muộn.\\
\bottomrule
\end{tabularx}
\end{table}

T38 giữ khóa xe, điểm đón và giờ đón sau xác nhận. Cấu hình $U$ bỏ cổng thời
gian đang so, nhưng vẫn giữ các khóa và tầng mục tiêu của hồ sơ vector cứng;
không đồng nhất nó với B1. Ở chế độ ghi nhận đón muộn, bộ chuyển đổi bỏ cận
trả neo theo cửa sổ đón đầu, vẫn giữ kiểm thời gian trên xe từ lúc đón thật.
Vì vậy hoàn tất trong T38 không chứng minh mọi cửa sổ vật lý ban đầu đều
được giữ. Mục~\ref{muc:bochon} xác định thứ tự mục tiêu tương ứng.

Lượt kiểm cấu hình đọc 20 đầu vào H6, 16 đầu vào phát triển tầng 2 và 49
đầu vào T38: cả 85 đều có đúng 96 nút, 9.120 cung và 108 yêu cầu. Băm kho
ZIP là \texttt{d9e86f33}, băm nguồn chuẩn hóa năm tệp là \texttt{34517a83};
cả hai được tính lại và khớp bản ghim. Băm DLL và cây công bố bộ thực thi
được kiểm riêng. Băm cây bộ chuyển đổi H6 v5 là trường đã ghi trong biên
nhận, có cả mã byte Python; khóa luận không nhận đã dựng lại độc lập cây
lịch sử ấy chỉ từ commit nguồn. Hồ sơ đầy đủ được giữ trong
\path{reproducibility-config.json} của gói kiểm và trong bản kê từng lần chạy.
